\documentclass[12pt]{article}
\usepackage[utf8]{inputenc}
\usepackage[T1]{fontenc}
\usepackage{hyperref}
\title{Uvbie Language Resources for K-12 Curriculum Development}
\author{Compiled by Claude Sonnet 5.5 (claude-sonnet-5-5) for Andr\'as Kornai}
\date{}
\begin{document}
\maketitle

\section{Introduction}
Uvbie (ISO 639-3 \texttt{evh}; also called Uvwie, Effurun, Evrie) is a Southwestern Edoid language of Delta State, Nigeria, with about 20,000 speakers. Sources disagree on whether it is separate from Urhobo; the Uvwie and Okpe clans regard their languages as distinct. Essentially no educational material specific to Uvbie was found.

\section{Primary Educational Resources}
None found. No primer, textbook or school programme for Uvbie appeared in the searches.

\section{Linguistic and Reference Materials}
\begin{itemize}
\item Uvbie language (Wikipedia): \url{https://en.wikipedia.org/wiki/Uvbie_language}
\item Southwestern Edoid languages (Wikipedia): \url{https://en.wikipedia.org/wiki/Southwestern_Edoid_languages}
\item Wikitongues language page: \url{https://wikitongues.org/languages/evh/}
\item Uvwie, the people and area (Wikipedia): \url{https://en.wikipedia.org/wiki/Uvwie}
\end{itemize}

\section{Religious and Community Resources}
Joshua Project reports Bible translation as started, with limited local-language Scripture: \url{https://www.joshuaproject.net/people_groups/15753/NI}

\section{Digital and Audio Resources}
None found beyond the Wikitongues page above.

\section{Curriculum Development Context}
The language may be shifting toward Warri Pidgin according to one source. Any curriculum would have to be built from scratch with local community input, possibly drawing on closely related Urhobo materials, which were not examined here.
\end{document}